% ═══════════════════════════════════════════════════════════════════════════════
%  CAPÍTULO — Proxy
% ═══════════════════════════════════════════════════════════════════════════════

\chapter{Proxy}

% ─── Situación Problema ──────────────────────────────────────────────────────
\section{Situación Problema}

Airbnb almacena información sensible de los anfitriones: número de
teléfono personal, dirección exacta del inmueble y datos de contacto
directo. Esta información solo debe ser accesible para el huésped una
vez que su reserva ha sido \emph{confirmada y pagada}; hasta ese momento,
la ficha de la propiedad muestra únicamente datos aproximados ---el
barrio y la ciudad, pero nunca la dirección exacta ni el teléfono real---.

El problema no es solo de permisos: el módulo que renderiza la ficha
de una propiedad no debería tener que conocer las reglas de negocio
que determinan qué información mostrar según el estado de la reserva.
Si esa lógica recae en la capa de presentación, se repite en cada punto
donde se muestran datos del anfitrión, es fácil de omitir y mezcla
responsabilidades de visualización con decisiones de acceso.

% ─── Solución ────────────────────────────────────────────────────────────────
\section{Solución}

El patrón Proxy introduce \texttt{ProxyPerfilAnfitrion}, que implementa
exactamente la misma interfaz \texttt{PerfilAnfitrion} que el objeto
real. El módulo de presentación solo conoce esa interfaz y solicita
los datos sin ninguna condición adicional. Internamente, el proxy
intercepta cada solicitud, consulta el estado de la reserva y decide
de forma autónoma qué devolver: si la reserva está confirmada, delega
al objeto real y retorna los datos completos; en caso contrario, retorna
una versión censurada donde el teléfono aparece enmascarado y la
dirección se reduce al nombre del barrio.

El cliente recibe siempre una respuesta a través de la misma interfaz
y en ningún momento toma ni conoce la decisión de qué versión de los
datos está viendo. Toda la política de acceso reside exclusivamente en
el proxy, lo que garantiza que no pueda ser ignorada por descuido,
centraliza su mantenimiento en un único lugar y deja al servicio real
completamente libre de lógica de autorización.

% ─── Diagrama de clases ─────────────────────────────────────────────────────
\begin{figure}[H]
    \centering
    % \includegraphics[width=\textwidth]{imagenes/abstract_factory_diagrama.png}
    \caption{Diagrama de clases del patrón Abstract Factory}
\end{figure}


% ─── Roles de las Clases ────────────────────────────────────────────────────
\section{Roles de las Clases}

% Explicar la responsabilidad de cada clase/interfaz

\begin{itemize}
    \item \textbf{AbstractFactory:}
    \item \textbf{ConcreteFactory:}
    \item \textbf{AbstractProduct:}
    \item \textbf{ConcreteProduct:}
    \item \textbf{Client:}
\end{itemize}


% ─── Main ───────────────────────────────────────────────────────────────────
\section{Main}

% Explicar qué realiza el programa principal

\begin{lstlisting}[language=Java, caption=NombreClase]
 
\end{lstlisting}


% ─── Salida del Main ────────────────────────────────────────────────────────
\section{Salida del Main}

\begin{lstlisting}[language=Java, caption=NombreClase]
 
\end{lstlisting}


% ─── Clases ─────────────────────────────────────────────────────────────────
\section{Clases}

% ─── Clase 1 ────────────────────────────────────────────────────────────────
\subsection{NombreClase}

\begin{lstlisting}[language=Java, caption=NombreClase]
 
\end{lstlisting}


% ─── Clase 2 ────────────────────────────────────────────────────────────────
\subsection{NombreClase}

\begin{lstlisting}[language=Java, caption=NombreClase]
 
\end{lstlisting}


% ─── Clase 3 ────────────────────────────────────────────────────────────────
\subsection{NombreClase}

\begin{lstlisting}[language=Java, caption=NombreClase]
 
\end{lstlisting}


% ─── Conclusiones ───────────────────────────────────────────────────────────
\section{Conclusiones}

% Explicar:
% - Ventajas del patrón
% - Cuándo usarlo
% - Beneficios obtenidos
% - Posibles desventajas
% - Relación con principios SOLID